% !TEX program = pdflatex
%
% DiSE — Adaptive Symbolic Importance Partitioning for Distribution-Aware
% Reliability Estimation
%
% Target venue: International Symposium on Automated Technology for
% Verification and Analysis (ATVA).
% Template: Springer LNCS (llncs).
%
% Build: pdflatex main && bibtex main && pdflatex main && pdflatex main
% Or:    latexmk -pdf main

\documentclass[runningheads]{llncs}

% ---------------------------------------------------------------------------
% Packages
% ---------------------------------------------------------------------------
\usepackage[T1]{fontenc}
\usepackage[utf8]{inputenc}
\usepackage{amsmath, amssymb}
% NB: amsthm conflicts with llncs's own theorem environments; we use
% llncs's \spnewtheorem and the built-in proof.
\usepackage{mathtools}
\usepackage{booktabs}
\usepackage{graphicx}
\usepackage{xcolor}
\usepackage{hyperref}
% cleveref / enumitem / algorithm / algpseudocode are nice-to-have but
% not in every minimal TeX distribution. To keep this draft buildable
% without extra installs, we use \providecommand fallbacks and the
% plain LaTeX enumerate/itemize. On a full TeX Live install, uncomment
% the lines below for typed cross-references and Algorithm boxes.
% \usepackage{cleveref}
% \usepackage{enumitem}
% \usepackage{algorithm}
% \usepackage{algpseudocode}
\providecommand{\Cref}[1]{Section~\ref{#1}}
\providecommand{\cref}[1]{section~\ref{#1}}
\usepackage{listings}
% microtype is nice but needs scalable fonts; disable for the minimal
% build. Comment out the next line if you want microtype on a full
% TeX install.
% \usepackage{microtype}

% ---------------------------------------------------------------------------
% Theorem environments (llncs already defines some)
% ---------------------------------------------------------------------------
\spnewtheorem*{assumption*}{Assumption}{\bfseries}{\itshape}

% ---------------------------------------------------------------------------
% Notation macros — keep this list short and consistent.
% ---------------------------------------------------------------------------
\newcommand{\inX}{\mathcal{X}}
\newcommand{\inY}{\mathcal{Y}}
\newcommand{\dist}{\mathcal{D}}
\newcommand{\prog}{P}
\newcommand{\prop}{\varphi}
\newcommand{\mutrue}{\mu}
\newcommand{\muhat}{\hat{\mu}}
\newcommand{\region}{R}
\newcommand{\partset}{\Pi}
\newcommand{\openleaves}{\partset^{\mathrm{open}}}
\newcommand{\closedleaves}{\partset^{\mathrm{closed}}}
\newcommand{\Wopen}{W_{\mathrm{open}}}
\newcommand{\epsstat}{\varepsilon_{\mathrm{stat}}}
\newcommand{\hwhat}{\hat{w}}
\newcommand{\hpathat}{\hat{p}}
\newcommand{\smt}{\textsc{Smt}}
\newcommand{\unsat}{\textsf{unsat}}
\newcommand{\sat}{\textsf{sat}}
\newcommand{\unknown}{\textsf{unknown}}
\newcommand{\toolname}{\textsc{DiSE}}
\newcommand{\algname}{\textsc{Asip}}
\newcommand{\code}[1]{\texttt{#1}}

% Listings styling (we mostly avoid raw code in the body but keep this for
% the implementation section).
\lstset{
  basicstyle=\ttfamily\small,
  breaklines=true,
  showstringspaces=false,
  columns=fullflexible,
}

% ---------------------------------------------------------------------------
% Document
% ---------------------------------------------------------------------------
\begin{document}

\title{Measuring the Frequency of Errors}
\titlerunning{Measuring the Frequency of Errors}

% \author{Anonymous}
\author{Aalok Thakkar\orcidID{0000-0000-0000-0000}}
\authorrunning{A. Thakkar}
\institute{Affiliation\\
  \email{thakkarsciencestrategy@gmail.com}}

\maketitle

\begin{abstract}
  We study \emph{distribution-aware reliability estimation}: given a
  deterministic program $\prog$, a discrete distribution $\dist$ over
  its inputs, and a Boolean property $\prop$, certify a confidence
  interval for $\Pr_{x \sim \dist}[\prop(\prog(x)) = 1]$ using as few
  program evaluations as possible.  Rare-event programs defeat plain
  sampling, while exact symbolic techniques fail on programs with
  non-linear arithmetic, unbounded loops, or non-uniform input
  distributions.

  We present \toolname{}, an adaptive partitioning algorithm that
  bridges the two regimes.  \toolname{} maintains a frontier of
  path-condition regions and at each step selects between
  \textsc{Sample} and SMT-driven \textsc{Refine} under a unified
  gain-per-cost criterion.  The certified half-width decomposes as
  $\Pr\!\bigl[|\muhat - \mutrue| \le \epsstat + \Wopen\bigr] \ge
  1 - \delta$, where $\epsstat$ aggregates per-region statistical
  uncertainty and $\Wopen$ is the total probability mass of regions
  not yet symbolically certified.  This decomposition is
  simultaneously the algorithm's correctness theorem and its
  scheduler: \textsc{Sample} reduces $\epsstat$, \textsc{Refine}
  reduces $\Wopen$, and the algorithm terminates as soon as their sum
  drops below the target $\varepsilon$.

  The per-region statistical bound instantiates the predictable-plug-in
  empirical-Bernstein confidence sequence of Waudby-Smith and
  Ramdas---closed-form, variance-adaptive, and anytime-valid under
  \toolname{}'s data-dependent stopping rule.  The closure rule is
  single-branch: a region is certified if and only if an SMT oracle
  proves path-determinism, with no sample-based fallback, yielding
  a one-lemma soundness contract.

  We implement \toolname{} as a Python prototype with a Z3 backend
  and evaluate it on twelve integer benchmarks drawn from established
  sources (Hacker's Delight, CLRS, the CERT~C Secure Coding standard,
  the Collatz literature) plus a parametric rare-event program.  Two
  ablations decompose the contribution of the bound and of the
  stratification: holding the bound constant, \toolname{}'s symbolic
  stratification reaches the same half-width with $9\times$ fewer
  samples; holding the adaptive-allocation strategy constant,
  symbolic guidance is $4\times$ tighter than variance-driven random
  stratification.  On a parametric rare-event program, \toolname{}
  certifies $\mutrue$ to half-width exactly zero in $\Theta(1)$
  samples in the sweet-spot regime $\mutrue \in [5\times 10^{-4},
  2\times 10^{-3}]$---a regime where plain Monte Carlo's
  $\Theta(1/(\mutrue\varepsilon^2))$ sample complexity is prohibitive.
\end{abstract}

\keywords{Probabilistic program analysis \and Statistical model
  checking \and Reliability estimation \and Symbolic execution \and
  Anytime-valid confidence intervals.}

\section{Introduction}\label{sec:intro}

Consider a network packet-routing function deployed on a server
handling millions of requests per day.  The function computes a hash
of each packet's header and maps it to one of $k$ queues; a collision
sends two packets to the same queue, degrading throughput.  The
production traffic distribution is far from uniform---certain header
patterns are orders of magnitude more common than others.  The
engineering team knows that edge-case collisions exist, because a
fuzzer found one.  What they need to know is the \emph{rate}: what
fraction of production inputs actually trigger the defect?  A single
counterexample does not answer this question.

This scenario is an instance of \emph{distribution-aware reliability
  estimation}.  Given a deterministic program $\prog : \inX \to \inY$,
a discrete distribution $\dist$ over $\inX$, and a Boolean property
$\prop : \inY \to \{0,1\}$, the \emph{operational reliability} of
$\prog$ under $\dist$ with respect to $\prop$ is
\begin{equation}\label{eq:mu-def}
  \mutrue \;=\; \Pr_{x \sim \dist}\!\big[\prop(\prog(x)) = 1\big].
\end{equation}
Reliability in this sense arises in \emph{operational testing}, where
a developer must certify that a production input distribution produces
a failure rate below a contractual threshold; in \emph{probabilistic
  program
  verification}~\cite{geldenhuys2012pse,filieri2015tacas}, where the
question is inherently quantitative; and in \emph{distribution-aware
  property-based
  testing}~\cite{maciver2019hypothesis,claessen2000quickcheck}, where
quantitative coverage supplements counterexample search.

\paragraph{The challenge.}
Three difficulties set reliability estimation apart from conventional
program analysis.

\emph{The answer is a real number, not a Boolean.}  Standard
verification asks $\forall x\colon \prop(\prog(x))$.  Reliability asks
for the \emph{measure} of the failure set $\{x : \prop(\prog(x))=1\}$
under $\dist$.  Counterexample-driven techniques produce a single
witness; bounding the probability of the entire set requires a
different methodology.

\emph{The rare-event regime matters most and is hardest.}  When
$\mutrue \ll 1$, plain Monte Carlo requires $\Theta(1/(\mutrue
\varepsilon^2))$ samples to certify a half-width of $\varepsilon$---a
prohibitive budget when $\mutrue = 10^{-4}$ and $\varepsilon =
10^{-5}$.  Yet rare events are precisely where operational reliability
matters: a defect that occurs once per million inputs under a uniform
distribution may occur once per hundred inputs under production
traffic.

\emph{The certificate is the deliverable.}  A point estimate for
$\mutrue$ without an accompanying confidence interval is operationally
useless: one cannot determine whether $\hat{\mu} = 0.003$ is
consistent with a contractual bound of $\le 0.001$ unless the
estimator's uncertainty is quantified.  Sound, tight intervals---not
just point estimates---are the engineering artifact that reliability
estimation must produce.

\paragraph{Existing approaches.}
Two families of techniques address this problem, each with
complementary strengths and limitations.

\emph{Probabilistic symbolic execution}
(PSE)~\cite{geldenhuys2012pse,filieri2015tacas,borges2014compositional}
enumerates program paths, encodes each path condition as a formula,
and computes its exact probability via model
counting~\cite{chakraborty2013approxmc,yang2023approxmc6}.  PSE is
exact where it applies.  However, path explosion forces incompleteness
on programs with unbounded loops or non-linear arithmetic, and
non-uniform input distributions require ad-hoc extensions of the
model-counting backend that are not always available.

\emph{Plain sampling}~\cite{sampson2014passert,kwiatkowska2011prism}
draws iid inputs, executes the program, and applies a concentration
inequality to the empirical failure rate.  Sampling is robust to all
sources of PSE difficulty, but inherits the $\Theta(1/\varepsilon^2)$
rate of independent estimation.  In the rare-event regime, the sample
budget required to reach a useful half-width is simply too large.

\paragraph{Our approach.}
We propose a middle course between exact symbolic enumeration and
unguided sampling.  \toolname{} maintains a \emph{frontier}: a tree whose leaves
partition the input domain $\inX$ into path-condition regions.  At
each step the algorithm chooses between two actions:
\textsc{Sample}$(\pi)$, which draws additional executions from a
region $\pi$ to tighten its statistical confidence bound, and
\textsc{Refine}$(\pi, b)$, which invokes an SMT solver to split $\pi$
along a branch predicate $b$ harvested from observed executions,
potentially certifying that an entire child region is all-pass or
all-fail under $\prop$.  The choice is governed by a single
gain-per-cost criterion applied uniformly to both action types.

The key insight is a two-part decomposition of the certified
half-width:
\begin{equation}\label{eq:hw-decomp}
  \Pr\!\big[\,|\muhat - \mutrue| \;\le\; \epsstat + \Wopen\,\big]
  \;\ge\; 1 - \delta,
\end{equation}
where $\epsstat$ aggregates per-region statistical uncertainty and
$\Wopen$ is the total probability mass of regions not yet
symbolically certified.  Every \textsc{Sample} action reduces
$\epsstat$; every successful \textsc{Refine} reduces $\Wopen$.  The
algorithm terminates as soon as $\epsstat + \Wopen \le \varepsilon$.
There are no phases, no hand-tuned switching heuristics: the
decomposition~\eqref{eq:hw-decomp} is simultaneously the correctness
argument and the scheduler.

Two design choices distinguish \toolname{} from prior work.  First,
the per-region statistical bound instantiates the predictable-plug-in
empirical-Bernstein (PrPl-EB) confidence sequence of Waudby-Smith and
Ramdas~\cite{waudbysmith2024betting}, which is closed-form,
variance-adaptive, and \emph{anytime-valid}: it remains sound under
\toolname{}'s data-dependent stopping rule, strictly tightening the
union-bounded Wilson construction used in earlier statistical model
checking
tools~\cite{boyer2013plasma,hartmanns2014modest}.  Second, the
closure rule is single-branch: a region is certified if and only if
the SMT backend proves path-determinism ($F_\pi \wedge
\neg\,\text{path}_\pi$ is \unsat{}).  There is no sample-based
fallback.  The price is wider intervals on programs whose closure is
symbolically intractable; the gain is a one-lemma soundness contract
and the elimination of a class of silent bias that affects prior
prototypes.

\paragraph{Contributions.}
This paper makes three contributions.

\begin{enumerate}
\item \emph{The \toolname{} algorithm and its correctness proof.}  We
  introduce \toolname{} and prove that its reported interval $[L,U]$
  covers $\mutrue$ with probability at least $1-\delta$
  (\Cref{thm:hw-decomp}).  The proof reduces to two independently
  checkable invariants---mass-conservative refinement
  (\Cref{lem:mass-conservation}) and SMT-certified closure
  (\Cref{lem:closure-soundness})---plus an application of an
  anytime-valid concentration inequality per leaf.

\item \emph{Anytime-valid certification via PrPl-EB.}  We give the
  first application of the predictable-plug-in empirical-Bernstein
  confidence sequence in the context of program reliability
  estimation.  The bound is closed-form, variance-adaptive, and
  requires no prior knowledge of per-leaf variance, strictly
  tightening standard Hoeffding and Wilson constructions under
  adaptive stopping.

\item \emph{Implementation and evaluation.}  We implement \toolname{}
  as \toolname{}, a Python prototype with a Z3 backend and a benchmark
  suite of $12$ integer programs.  On rare-event benchmarks whose path
  structure aligns with the input distribution, \toolname{} reduces
  sample counts by up to two orders of magnitude relative to plain and
  stratified Monte Carlo.
\end{enumerate}

\paragraph{Paper organisation.}
\Cref{sec:problem} formally states the problem and defines the input
distribution class.  \Cref{sec:related} surveys related work.
\Cref{sec:algorithm} presents \toolname{} and proves
\Cref{thm:hw-decomp}.  \Cref{sec:implementation} describes the
\toolname{} implementation.  \Cref{sec:results} reports experimental
results.  \Cref{sec:conclusion} concludes.

\section{Problem Statement}\label{sec:problem}

Fix a probability space $(\Omega, \mathcal{F}, \mathbb{P})$ and a
finite input domain $\inX \subseteq \mathbb{Z}^d$.  Throughout the
paper we make the following assumption on the program.

\begin{assumption*}\label{ass:program}
  The program $\prog : \inX \to \inY$ is \emph{deterministic}: the
  output $\prog(x)$ is determined entirely by the input $x$, with no
  dependence on any internal source of randomness.  Programs that use
  randomness must be wrapped to expose the random seed as an explicit
  input.
\end{assumption*}

We consider input distributions drawn from the following class.

\begin{definition}[Product-form discrete distributions, class (D1)]%
  \label{def:d1}
  A distribution $\dist$ over $\inX = \mathbb{Z}^d$ is in class
  \emph{(D1)} if it factors as
  $\dist(x_1, \ldots, x_d) = \prod_{i=1}^d \dist_i(x_i)$,
  where each marginal $\dist_i$ has a closed-form cumulative
  probability-mass function over integer intervals:
  $\dist_i(\{a, \ldots, b\}) = \dist_i(b) - \dist_i(a-1)$.
  Uniform, bounded geometric, categorical, truncated Poisson, and
  finite Bernoulli mixtures all belong to (D1).
\end{definition}

The (D1) restriction is the source of \toolname{}'s central
variance-reduction mechanism.  For any axis-aligned box
$R = \prod_i [a_i, b_i]$, the probability
$\Pr_\dist[X \in R] = \prod_i \dist_i(\{a_i, \ldots, b_i\})$ is
available in closed form and carries \emph{zero estimation variance}.
Bayes-net-structured and general discrete distributions fall outside
this class and are left to future work (\Cref{sec:conclusion}).

Let $\prop : \inY \to \{0,1\}$ be a Boolean output property.  The
operational reliability $\mutrue \in [0,1]$ is defined as
in~\eqref{eq:mu-def}.

\paragraph{Problem definition.}
Given $\prog$, $\dist \in \text{(D1)}$, $\prop$, a target half-width
$\varepsilon \in (0,1]$, and a confidence parameter $\delta \in (0,1)$,
produce an estimator $\muhat$ and a two-sided interval
$[L,U] \subseteq [0,1]$ satisfying the \emph{soundness contract}
\begin{equation}\label{eq:soundness-contract}
  \mathbb{P}\!\big[\,L \le \mutrue \le U\,\big] \;\ge\; 1 - \delta,
  \qquad
  U - L \;\le\; 2\varepsilon,
\end{equation}
using as few evaluations of $\prog$ as possible.  The probability
in~\eqref{eq:soundness-contract} is over the algorithm's internal
randomness and the iid samples drawn from $\dist$; the true
reliability $\mutrue$ is a fixed constant.

\paragraph{Assertion-violation framing.}
The output-property formulation~\eqref{eq:mu-def} subsumes the
classical assertion-violation framing.  For a program annotated with
an \code{assert} statement, the failure probability
$\Pr_\dist[\prog\text{ raises }\mathsf{AssertionError}]$ is recovered
by wrapping $\prog$ in a handler that converts the targeted exception
into a Boolean failure flag and setting $\prop(y) = [y = 1]$.
Worked examples of both framings appear in \Cref{sec:results}.

\paragraph{Why the problem is hard.}
\toolname{} addresses three inherent difficulties, each of which is
mitigated by a distinct component of its design.
\begin{enumerate}
\item \emph{Real-valued answer.}  Reliability is a measure, not a
  predicate.  \toolname{} maintains a fractional estimator over a
  partition of $\inX$, accumulating evidence from both sampling and
  symbolic reasoning.
\item \emph{Rare events.}  When $\mutrue \approx 0$ or $\mutrue
  \approx 1$, plain Monte Carlo is sample-inefficient.  \toolname{}
  stratifies the input domain by path condition, concentrating samples
  in regions of high variance and using closed-form mass computations
  in certified regions.
\item \emph{Sound, tight intervals.}  A point estimate without a
  certificate is operationally useless.  \toolname{} carries an
  explicit anytime-valid confidence bound at every step, so the
  interval it reports is sound regardless of when it stops.
\end{enumerate}

\section{Related Work}\label{sec:related}

\algname{} sits at the intersection of five research lines.  We
position it against each.

\paragraph{Probabilistic symbolic execution.}
Geldenhuys, Dwyer, and Visser~\cite{geldenhuys2012pse} introduced
probabilistic symbolic execution (PSE), combining symbolic execution
with model counting (LattE, Barvinok~\cite{koppe2008barvinok}) to
compute exact path probabilities for affine path conditions over
uniform inputs.  Filieri et al.~\cite{filieri2015tacas} extended PSE
to structured data; Borges et al.~\cite{borges2014compositional}
introduced compositional model counting.  Subsequent work, surveyed in
Pasareanu's synthesis lectures~\cite{pasareanu2022lecture}, has
broadened the class of treatable path conditions and input
distributions.

PSE is exact where it applies; \algname{} is statistical and
produces $(1-\delta)$-coverage intervals.  On programs where PSE
applies---bounded linear-integer-arithmetic paths and tractable input
distributions---PSE strictly dominates.  \algname{} fills the gap
when programs have unbounded loops, non-linear arithmetic, or
non-uniform input distributions.  When every region in \algname{}'s
frontier admits a closed-form mass and the closure rule succeeds on
every leaf, the two approaches agree exactly.

\paragraph{Statistical and probabilistic model checking.}
PRISM~\cite{kwiatkowska2011prism} and Storm~\cite{hensel2022storm}
verify quantitative temporal properties on Markov chains and MDPs.
Statistical variants---PLASMA-Lab~\cite{boyer2013plasma}, the Modest
Toolset~\cite{hartmanns2014modest}, and Storm's SMC
front-end~\cite{schoolderman2024storm}---discharge the same
specifications by sampling the underlying model.  These tools operate
on dedicated modelling languages (PRISM, JANI, Modest); \algname{}
operates directly on Python source, treating the program itself as the
model.  The trade-off is complementary: \algname{} cannot reason about
nondeterminism or scheduler choices, while PMC tools cannot reason
about source-level loops without a hand-built model.

\paragraph{Probabilistic assertions.}
Sampson et al.~\cite{sampson2014passert} introduced
\emph{passerts}---probabilistic assertions in C programs verified by
sampling with Hoeffding or Bayesian bounds.  \toolname{}'s
\code{failure\_probability} API targets the same query class.  The
difference is the bound and the sampling strategy: passerts use plain
Hoeffding sampling over the whole input space, whereas \algname{}
stratifies by symbolic path condition and uses a variance-adaptive
anytime-valid bound~\cite{waudbysmith2024betting}.  On programs whose
path structure concentrates probability mass, \algname{} reaches the
same certified half-width with substantially fewer samples; on
programs with flat, well-mixed distributions the advantage narrows.

\paragraph{Model counting.}
Approximate model counting~\cite{chakraborty2013approxmc,yang2023approxmc6}
provides PAC-$(\varepsilon,\delta)$ counts for propositional
satisfiability problems.  Recent work has added formal certificates
in Isabelle~\cite{kiesl2024certified}.  \algname{} does not currently
invoke a model counter---it uses importance sampling to estimate mass
on non-axis-aligned regions---but a natural extension would delegate
the region-mass subproblem to a certified counter, retaining
$(1-\delta)$-coverage with tighter constants.

\paragraph{Anytime-valid concentration inequalities.}
Classical Wilson and Hoeffding bounds are valid only at a
\emph{fixed} sample size; applying them under data-dependent stopping
requires a union bound over sample counts.  The mixture-martingale
literature, culminating in Howard et al.~\cite{howard2021time},
provides time-uniform Chernoff bounds via non-negative
supermartingales.  Waudby-Smith and
Ramdas~\cite{waudbysmith2024betting} sharpen this with a closed-form
predictable-plug-in empirical-Bernstein (PrPl-EB) confidence sequence
that is variance-adaptive and matches the asymptotic width of
Bennett's inequality without requiring the variance to be known.
\algname{} uses PrPl-EB per leaf with Bonferroni correction over the
(bounded) leaf count.

\paragraph{Property-based testing.}
QuickCheck~\cite{claessen2000quickcheck} and
Hypothesis~\cite{maciver2019hypothesis} search for a single input that
violates a property---a different goal from bounding the failure rate.
We provide a thin adapter (\code{dise.integrations.hypothesis}) that
converts Hypothesis strategies into (D1) distributions, so the same
property test can drive both edge-case search and operational
reliability estimation.

\paragraph{Summary.}
\algname{}'s technical contribution lies between the exact PSE line,
which it does not displace where PSE applies, and the sampling-only
literature, which it improves upon via symbolic stratification.  No
prior algorithm combines SMT-driven adaptive partitioning, closed-form
mass computation on axis-aligned (D1) regions, and an anytime-valid
empirical-Bernstein bound in a single framework.

\section{The \algname{} Algorithm}\label{sec:algorithm}

We describe \algname{} (\emph{Adaptive Symbolic Importance
  Partitioning}).  \Cref{sec:frontier} introduces the frontier data
structure.  \Cref{sec:hw-decomp} states the central decomposition
theorem, which is the load-bearing correctness result of the paper.
\Cref{sec:asip-loop} gives the driver loop.

\subsection{The Frontier}\label{sec:frontier}

\algname{} maintains a \emph{frontier}: a tree of path-condition
regions $\{R_\pi\}_{\pi \in \partset}$ whose leaves form a partition
of $\inX$.  Each leaf $\pi$ carries four pieces of data:
\begin{itemize}
\item a symbolic region formula $F_\pi$ in the theory of
  linear integer/bitvector arithmetic extended with equality and
  propositional connectives;
\item a probability estimate $\hwhat_\pi \approx \Pr_\dist[X \in R_\pi]$;
\item a sample-based hit-rate estimate
  $\hpathat_\pi \approx \Pr[\prop(\prog(X))=1 \mid X \in R_\pi]$;
\item a lifecycle status in
  $\{\mathsf{OPEN}, \mathsf{CLOSED\text{-}T},
    \mathsf{CLOSED\text{-}F}, \mathsf{EMPTY}\}$.
\end{itemize}

A leaf is \textsc{open} if its hit rate is not yet symbolically
certified.  It is \textsc{closed-t} or \textsc{closed-f} if the SMT
backend has proved that $\prop(\prog(x)) = 1$ or $\prop(\prog(x)) = 0$
holds uniformly on the region.  It is \textsc{empty} if SMT has shown
$F_\pi \equiv \bot$.  The overall reliability estimator is
\begin{equation}\label{eq:muhat-def}
  \muhat \;=\; \sum_{\pi \in \partset} \hwhat_\pi \cdot \hpathat_\pi.
\end{equation}

\begin{lemma}[Mass-conservative refinement]\label{lem:mass-conservation}
  When \algname{} refines a leaf $\pi$ on a clause $b$, producing
  children $\pi{\wedge}b$ and $\pi{\wedge}\neg b$, the children's
  estimated masses satisfy
  $\hwhat_{\pi \wedge b} + \hwhat_{\pi \wedge \neg b} = \hwhat_\pi$
  exactly (modulo importance-sampling noise on non-axis-aligned
  regions, which is bounded in \Cref{sec:implementation}).
\end{lemma}

\begin{proof}[Proof sketch]
  For axis-aligned LIA children, the closed-form mass of each child
  is a product of marginal-PMF differences over disjoint coordinate
  ranges, which sum exactly to the parent's mass.  For general
  regions, \algname{} draws one importance-sampling batch from
  $\dist \mid R_\pi$, partitions it by $b$, and splits $\hwhat_\pi$
  proportionally: $\hwhat_{\pi \wedge b} = \hwhat_\pi \cdot
  \hpathat_{b \mid \pi}$ and symmetrically for $\neg b$.  The split
  conserves $\hwhat_\pi$ by construction.
\end{proof}

\begin{lemma}[Closure soundness]\label{lem:closure-soundness}
  Suppose the SMT oracle returns \unsat{} for
  $F_\pi \wedge \neg\,\mathrm{path}_\pi$, where
  $\mathrm{path}_\pi$ is the conjunction of branch clauses observed
  at $\pi$.  Then every $x \in R_\pi$ traverses
  $\mathrm{path}_\pi$ in $\prog$.  If additionally every
  observation at $\pi$ agrees on the same $\prop$-value
  $v \in \{0,1\}$, then $\prop(\prog(x)) = v$ for every
  $x \in R_\pi$.
\end{lemma}

\begin{proof}[Proof sketch]
  Suppose some $x \in R_\pi$ takes a different path.  Then some
  branch clause $b_j$ of $\mathrm{path}_\pi$ is violated by $x$,
  so $x \models F_\pi \wedge \neg\,\mathrm{path}_\pi$, contradicting
  the oracle's \unsat{} verdict (by soundness of the SMT backend).
  Determinism of $\prog$ then forces $\prop(\prog(x)) = v$ uniformly
  on $R_\pi$.
\end{proof}

\subsection{The Central Decomposition Theorem}\label{sec:hw-decomp}

Let $\Pi_T$ denote the frontier at the algorithm's stopping time $T$,
decomposed as
$\Pi_T = \openleaves \cup \closedleaves \cup \Pi^{\mathsf{EMPTY}}_T$.
Define the two half-width terms:
\begin{align}
  \Wopen(T) \;&=\; \sum_{\pi \in \openleaves} \hwhat_\pi,\\
  \epsstat(T) \;&=\; \sum_{\pi \in \openleaves}
    \hwhat_\pi \cdot W_\pi(T),
\end{align}
where $W_\pi(T)$ is an anytime-valid $(1-\delta_\pi)$-half-width for
the per-leaf hit rate, with Bonferroni allocation
$\delta_\pi = \delta / K_T$ over the $K_T \le K_{\max}$ leaves
created by the algorithm.

\begin{theorem}[Central decomposition; coverage at $1-\delta$]%
  \label{thm:hw-decomp}
  Under Assumption~1, Lemma~\ref{lem:mass-conservation},
  Lemma~\ref{lem:closure-soundness}, soundness of the SMT oracle, and
  an anytime-valid per-leaf bound $W_\pi$,
  \begin{equation}
    \mathbb{P}\!\Big[\,\big|\muhat_T - \mutrue\big|
      \;\le\; \epsstat(T) + \Wopen(T)\,\Big]
    \;\ge\; 1 - \delta.
  \end{equation}
\end{theorem}

\begin{proof}[Proof sketch]
  Decompose the estimation error:
  \[
    \muhat_T - \mutrue
    = \sum_{\pi \in \openleaves}
        \!\bigl(\hwhat_\pi\hpathat_\pi - w_\pi p_\pi\bigr)
      + \sum_{\pi \in \closedleaves}
        \!\bigl(\hwhat_\pi - w_\pi\bigr)p_\pi.
  \]
  The closed-leaf term vanishes for axis-aligned closed leaves
  (by the closed-form mass) and is bounded by
  Lemma~\ref{lem:mass-conservation} otherwise.  Each open-leaf term
  splits into a deterministic part---bounded in absolute value by
  $\hwhat_\pi$, summing to $\Wopen(T)$---and a stochastic part
  bounded with probability $\ge 1-\delta_\pi$ by
  $\hwhat_\pi W_\pi(T)$.  A union bound over the $K_T$ leaves
  controls the total stochastic deviation by $\epsstat(T)$ at
  confidence $1-\delta$.  The anytime property of $W_\pi$ makes
  the union bound valid at the data-dependent stopping time $T$.
\end{proof}

Theorem~\ref{thm:hw-decomp} is the load-bearing result.  Every
implementation choice in \Cref{sec:implementation} is justified by
reference to it, and every action the algorithm takes reduces either
$\epsstat$ or $\Wopen$.  The reported interval is the two-sided
projection
\begin{equation}\label{eq:interval}
  [L, U] \;=\;
  \bigl[\max(0,\,\muhat_T - \epsstat(T) - \Wopen(T)),\;
        \min(1,\,\muhat_T + \epsstat(T) + \Wopen(T))\bigr].
\end{equation}

\subsection{The Driver Loop}\label{sec:asip-loop}

At each iteration, \algname{} considers two action families.
\begin{description}
\item[\textsc{Sample}$(\pi, b)$.]
  Draw $b$ additional iid samples from $\dist \mid R_\pi$, increment
  $n_\pi$ by $b$, update $\hpathat_\pi$ and $W_\pi$.
  Expected gain: $\Delta\epsstat \approx \hwhat_\pi\cdot
  (W_\pi(n_\pi) - W_\pi(n_\pi+b))$.
  Cost: $b$ program evaluations.

\item[\textsc{Refine}$(\pi, b)$.]
  Select a branch clause $b$ observed at $\pi$.  Split $\pi$ into
  $\pi{\wedge}b$ and $\pi{\wedge}\neg b$, distributing $\hwhat_\pi$
  per Lemma~\ref{lem:mass-conservation}, then immediately attempt
  closure on each child (Lemma~\ref{lem:closure-soundness}).
  Expected gain: $\Delta\Wopen \approx \hwhat_\pi \cdot
  \Pr[\text{at least one child closes}]$.
  Cost: one SMT call (split) plus two (closure checks).
\end{description}

The scheduler selects $\arg\max_a\,\mathrm{gain}(a)/\mathrm{cost}(a)$,
where cost is normalised to one concolic run per unit.  The algorithm
terminates when $\epsstat(T) + \Wopen(T) \le \varepsilon$, or when
one of three optional caps fires: sample budget exhausted, wall-clock
limit reached, or no action exceeds a minimum gain-per-cost floor.

\begin{figure}[t]
  \caption{\algname{} (Adaptive Symbolic Importance Partitioning)}
  \label{alg:asip}
  \hrule\vspace{0.4em}
  \textbf{Input:} program $\prog$, distribution $\dist$, property
    $\prop$, target $(\varepsilon, \delta)$, optional budget $B$.\\
  \textbf{Output:} estimator $\muhat$ and certified interval $[L,U]$.
  \begin{flushleft}
  \begin{tabular}{rl}
    1: & $\partset \gets \{\pi_0\}$ with $F_{\pi_0} = \top$,
         $\hwhat_{\pi_0} = 1$ \hspace{1.5em}// single \textsc{open}
         leaf at root\\
    2: & bootstrap $n_{\mathrm{boot}}$ samples from $\dist$ at
         $\pi_0$;\; call \textsc{TryClose}($\pi_0$)\\
    3: & \textbf{while} not \textsc{Terminated}() \textbf{do}\\
    4: & \quad $\mathcal{A} \gets
         \{\textsc{Sample}(\pi,b),\,\textsc{Refine}(\pi,b)
           : \pi \in \openleaves\}$\\
    5: & \quad $a^\star \gets
         \arg\max_{a \in \mathcal{A}}\,\mathrm{gain}(a)/\mathrm{cost}(a)$\\
    6: & \quad execute $a^\star$;\; update $\partset$\\
    7: & \quad call \textsc{TryClose}($\pi$) for every newly-affected
         leaf\\
    8: & \textbf{end while}\\
    9: & \textbf{return} $(\muhat,\,[L,U])$ per
         \eqref{eq:muhat-def} and~\eqref{eq:interval}\\
  \end{tabular}
  \end{flushleft}
  \vspace{0.2em}\hrule
\end{figure}

\algname{} is parametrised by three engineering choices: the
candidate-clause heuristic for \textsc{Refine}, the per-leaf bound
construction $W_\pi$, and the SMT backend.  Each choice is discussed
in \Cref{sec:implementation}.

\section{Implementation}\label{sec:implementation}

\toolname{} is an open-source Python prototype
of approximately $5{,}000$ lines of code with $290+$ tests.  This
section covers the three engineering choices that turn
\Cref{alg:asip} into a working artifact: the concolic front-end, the
per-leaf bound $W_\pi$, and the SMT abstraction layer.

\subsection{Concolic Execution and Branch Capture}\label{sec:concolic}

A concolic run~\cite{sen2005cute,godefroid2005dart} executes $\prog(x)$
on a concrete input $x$ while simultaneously tracking the symbolic
conditions on $x$ that drove each branch taken.  \toolname{}
implements this via a tagged-integer wrapper: arithmetic operations on
tagged values propagate their symbolic representations, and comparisons
emit a branch record $(\text{clause},\text{direction})$ before
delegating to the concrete operator.

Each run produces the path condition (the conjunction of branch clauses
along the path taken), the property value $\prop(\prog(x))$, and the
concrete input.  Property evaluation is intentionally \emph{outside}
the tracer: $\prop$ is evaluated on the concrete output, and the
resulting Boolean is appended to the run record.  This division of
labour keeps path conditions small, since $\prop$ is typically a thin
projection (\code{$y \le k$}, \code{popcount($y$) $\ge 0$}) that
should not inflate the symbolic representation.

\subsection{The Per-Leaf Bound $W_\pi$}\label{sec:per-leaf-bound}

Theorem~\ref{thm:hw-decomp} requires an anytime-valid
$(1-\delta_\pi)$-half-width for each open leaf's hit rate.  \toolname{}
exposes four choices via the \code{method} parameter.

\paragraph{Wilson (\code{method="wilson"}).}
The classical fixed-$n$ Wilson score interval~\cite{wilson1927}.  Tight
on Bernoulli leaves at a pre-committed sample count; \emph{not} sound
under data-dependent stopping.  Appropriate only when the budget is
fixed in advance.

\paragraph{Anytime Wilson (\code{method="anytime"}).}
The Wilson interval evaluated at confidence $1 - \delta_n$ with
$\delta_n = 6\delta_\pi/(\pi^2 n^2)$, union-bounded over all $n$ via
$\sum_{n=1}^\infty 6/(\pi^2 n^2) = 1$.  Sound under data-dependent
stopping; loose by a $\sqrt{\log n}$ factor and a $\pi^2/6$
inflation.

\paragraph{Predictable-plug-in empirical Bernstein (\code{method="betting"}).}
Theorem~2 of Waudby-Smith and Ramdas~\cite{waudbysmith2024betting}.
For a Bernoulli history $X_1,\ldots,X_n$ at leaf $\pi$, the half-width
is
\begin{equation}\label{eq:prpl-eb}
  W_\pi \;=\;
    \frac{\log(2/\delta_\pi) + \sum_{i=1}^n v_i\,\psi_E(\lambda_i)}
         {\sum_{i=1}^n \lambda_i},
  \qquad
  \psi_E(\lambda) = \tfrac{-\log(1-\lambda) - \lambda}{4},
\end{equation}
with $v_i = 4(X_i - \hat\mu_{i-1})^2$ and predictable bet
\begin{equation}
  \lambda_t \;=\; \min\!\Bigg(
    \sqrt{\frac{2\log(2/\delta_\pi)}
              {\hat\sigma_{t-1}^2 \cdot t\log(t+1)}},\;c\Bigg),
  \quad c \in (0,1),
\end{equation}
where $\hat\mu_t$ and $\hat\sigma_t^2$ are Wilson-regularised running
mean and variance.  The default truncation is $c = 1/2$, as
recommended in~\cite[\S3.2]{waudbysmith2024betting}.  PrPl-EB is
anytime-valid, closed-form, and variance-adaptive: its asymptotic
width is $\sqrt{2\sigma^2\log(2/\delta_\pi)}$, matching Bennett's
inequality without requiring known variance.  This is the recommended
default for certified results and is strictly tighter than
\code{anytime} on low-variance leaves.

\paragraph{Classical and Maurer--Pontil Bernstein.}
Provided as \code{bernstein} and \code{empirical-bernstein} for
baseline comparison only; both are fixed-$n$ bounds and not sound
under data-dependent stopping.

\subsection{The Closure Rule}\label{sec:closure}

A leaf is closed only when the SMT backend certifies
path-determinism:
\begin{verbatim}
def try_close(node):
    if n_samples < min_samples: return False
    if not all observed paths identical: return False
    if not all observed phi values identical: return False
    if is_satisfiable(F_pi AND NOT observed_path) != "unsat":
        return False
    node.status = CLOSED_T if phi == 1 else CLOSED_F
    return True
\end{verbatim}
There is no sample-based fallback.  A leaf whose closure the SMT
backend cannot certify (\sat{} or \unknown{}) remains open and
contributes to $\Wopen$.  This makes the soundness contract a single
lemma (Lemma~\ref{lem:closure-soundness}), at the cost of wider
intervals on programs whose path structure is symbolically
intractable.

\subsection{SMT Abstraction}\label{sec:smt}

The closure check, region-emptiness check, and refinement-clause
selection all go through an \code{SMTBackend} protocol with three
implementations: a Z3 backend~\cite{demoura2008z3} for production use,
a \code{MockBackend} restricted to axis-aligned LIA that permits the
test suite to run without external dependencies, and a
\code{CachedBackend} wrapper that memoises solver queries (typical
cache hit rate: $40$--$60\%$ on the benchmark suite).  All backends
guarantee soundness of \sat/\unsat{}; returning \unknown{} is always
admissible.

\subsection{Refinement-Clause Selection}\label{sec:clause-selection}

\toolname{}'s \textsc{Refine} action requires selecting a branch clause
$b$ from the finite set harvested from observed paths at $\pi$.
\toolname{} scores each candidate by expected mass balance---intuitively,
the clause that partitions $\hwhat_\pi$ closest to 50/50, maximising
the probability that at least one child closes.  The heuristic is
deliberately conservative; a principled alternative based on Craig
interpolation is identified as the primary next development direction
(\Cref{sec:conclusion}).

\subsection{Reproducibility}\label{sec:repro}

\toolname{}'s source code, benchmark suite, and reproduction scripts
are available at \code{https://github.com/<anon>/dise} (URL anonymised
for review).  The Docker image and \code{scripts/reproduce.sh}
regenerate all results in \Cref{sec:results} in approximately
$30$--$60$ minutes on commodity hardware.  Seed control is
centralised: each \code{(method, benchmark, seed)} triple is fully
deterministic.

\section{Experimental Evaluation}\label{sec:results}

We evaluate \toolname{} along three axes.
\Cref{sec:rare-event-scaling} presents the \emph{headline experiment}:
sample efficiency as a function of the rare-event mass~$\mutrue$, on
a parametric program whose structure isolates DiSE's
variance-reduction mechanism.
\Cref{sec:ablations} reports two \emph{ablations} that decompose the
overall win into the contribution of the statistical bound and the
contribution of symbolic stratification.
\Cref{sec:suite} reports the full benchmark suite at a shared
budget.
We close with a discussion of the regimes where \toolname{} is the
right tool and where it is not (\Cref{sec:discussion}).

All experiments use $\delta = 0.05$, three random seeds per cell, and
the Z3 backend.  Hardware: a laptop-class machine.  All raw rows
($720$ per-cell measurements for the suite, $105$ for the rare-event
sweep) and aggregate JSON summaries are in the supplementary
material; reproduction is one command,
\code{scripts/reproduce.sh}.

\subsection{Comparators}\label{sec:comparators}

We compare against seven sampling-only methods and three \toolname{}
variants (\Cref{tab:comparators}).  The two ablations require five of
the sampling-only methods to be in the suite simultaneously; without
them, the experiment cannot distinguish \emph{what the bound
contributes} from \emph{what the stratification contributes}.

\begin{table}[t]
  \centering
  \caption{The ten comparators evaluated in this study.  The two
    ablations isolate orthogonal sources of \toolname{}'s advantage.}
  \label{tab:comparators}
  \small
  \begin{tabular}{lll}
    \toprule
    Method & Per-sample bound & Stratification \\
    \midrule
    \code{plain\_mc}             & Wilson (Bernoulli-tight) & none \\
    \code{plain\_mc\_hoeffding}  & Hoeffding $\sqrt{\log(2/\delta)/2n}$ & none \\
    \code{plain\_mc\_eb}         & Maurer--Pontil EB~\cite{maurer2009empirical} & none \\
    \code{plain\_mc\_betting}    & WSR PrPl-EB~\cite{waudbysmith2024betting} & none \\
    \code{quasi\_mc\_sobol}      & Wilson on Sobol points & low-discrepancy \\
    \code{stratified\_random}    & per-bucket Wilson + Bonferroni & $16$ random hash buckets \\
    \code{adaptive\_stratified}  & per-bucket Wilson + Bonferroni & $16$ buckets + Neyman allocation \\
    \midrule
    \code{dise\_wilson}          & Wilson per leaf + Bonferroni  & SMT path-condition partition \\
    \code{dise\_anytime}         & union-bound-in-time Wilson    & SMT path-condition partition \\
    \code{dise\_betting}         & WSR PrPl-EB per leaf          & SMT path-condition partition \\
    \bottomrule
  \end{tabular}
\end{table}

The two pivotal pairs are:

\textbf{Ablation~A (the bound contribution).}
\code{plain\_mc\_betting} and \code{dise\_betting} share the same
PrPl-EB statistical bound; they differ only in whether symbolic
stratification is applied.  Their gap measures the contribution of
stratification alone.

\textbf{Ablation~B (the stratification contribution).}
\code{adaptive\_stratified} (two-pass Neyman allocation across $16$
random hash buckets, in the spirit of
Carpentier--Munos~\cite{carpentier2011}) is the strongest
\emph{pure-sampling} stratifier.  Its gap to \code{dise\_betting}
measures the contribution of symbolic guidance over variance-driven
adaptive sampling.

\subsection{Benchmarks}\label{sec:benchmarks}

Every benchmark in the suite is drawn from a recognised source.
\Cref{tab:benchmarks} lists the twelve programs and their provenance.

\begin{table}[t]
  \centering
  \caption{The benchmark suite, by provenance.  $\mutrue^\star$ is
    the Monte Carlo reference value used to verify empirical
    coverage.  All benchmarks share the property
    $\prop(\prog(x)) = 1$ except where noted; assertion benchmarks
    fail by raising \code{AssertionError}.}
  \label{tab:benchmarks}
  \small
  \begin{tabular}{llrl}
    \toprule
    Benchmark & Property & $\mutrue^\star$ & Source \\
    \midrule
    \multicolumn{4}{l}{\emph{Hacker's Delight 2e~\cite{warren2012hd}
      (bitvector kernels):}}\\
    \code{popcount\_w6}   & $\ge 0$ (tautology) & $1.000$ & §5--2 \\
    \code{parity\_w6}     & $\in \{0,1\}$ (tautology) & $1.000$ & §5--2 \\
    \code{log2\_w6}       & $\ge -1$ (tautology)       & $1.000$ & §5--3 \\
    \midrule
    \multicolumn{4}{l}{\emph{CLRS~\cite{cormen2009clrs} / Knuth
      TAoCP~\cite{knuth1998taocp} (classical algorithms):}}\\
    \code{gcd\_steps\_le\_5\_BG(0.1,100)}    & steps $\le 5$  & $0.989$ & TAoCP §4.5.2 \\
    \code{modpow\_fits\_in\_4b\_m=37}        & output $< 16$  & $0.62$  & CLRS Ch.~31 \\
    \code{integer\_sqrt\_correct\_U(1,1023)} & spec holds      & $1.000$ & CLRS \\
    \code{sieve\_primality\_U(2,200)}        & is prime        & $0.23$  & classical \\
    \code{miller\_rabin\_w=2\_BG(0.05,200)}  & accepts         & $0.91$  & \cite{rabin1980} \\
    \midrule
    \code{collatz\_le\_30\_BG(0.05,200)}     & steps $\le 30$  & $0.74$  & Lagarias~\cite{lagarias2010collatz} \\
    \code{sparse\_trie\_depth\_le\_3\_U(0,63)}   & depth $\le 3$ & $1.000$ & TAoCP §6.3 \\
    \code{assertion\_overflow\_mul\_w=8\_U(1,31)} & $a \cdot b$ overflows & $0.401$ & CERT C INT32-C \\
    \code{coin\_machine\_U(1,9999)}          & output $= 1$    & $0.0099$ & pedagogical \\
    \bottomrule
  \end{tabular}
\end{table}

\subsection{Headline Experiment: Rare-Event Scaling}\label{sec:rare-event-scaling}

The headline question is: as $\mutrue$ shrinks, how many samples does
each method need to certify $\mutrue$ to a fixed half-width
$\varepsilon$?  Theory makes two competing predictions.  Plain Monte
Carlo's sample complexity is $\Theta(1/(\mutrue \varepsilon^2))$;
\toolname{} on a partition all of whose leaves admit closed-form mass
should be $\Theta(1)$, independent of $\mutrue$, once the bootstrap
phase has observed an example from each region.

To isolate this claim, we introduce a parametric program with three
axis-aligned regions (\Cref{fig:rare-slice-prog}).  All path
conditions are linear over the single integer variable, so
\toolname{}'s refinement produces three axis-aligned leaves whose
mass is computed in closed form with \emph{zero} estimation variance.

\begin{figure}[t]
\centering
\fbox{\parbox{0.85\textwidth}{\small\ttfamily
def rare\_slice(x):\\
\hspace*{1em}if x \textless\ 1000:                          return 0   \# below the slice\\
\hspace*{1em}if x \textgreater= 1000 + rare\_width:         return 0   \# above the slice\\
\hspace*{1em}return 1                                                  \# inside the slice
}}
\caption{The parametric rare-event program.  Varying
  $\code{rare\_width}$ from $5000$ to $1$ moves
  $\mutrue = \code{rare\_width}/N$ across four orders of magnitude
  under $x \sim \mathrm{Uniform}(1, 9999)$.  All path conditions are
  axis-aligned, so \toolname{}'s refinement produces three
  closed-form leaves.}
\label{fig:rare-slice-prog}
\end{figure}

\begin{figure}[t]
  \centering
  \includegraphics[width=0.85\textwidth]{00_rare_event_scaling.pdf}
  \caption{Samples required to certify $\mutrue$ to half-width
    $\varepsilon = 0.005$ at $\delta = 0.05$, as a function of the
    true rare-event mass $\mutrue$, on the parametric program of
    \Cref{fig:rare-slice-prog}.  Three seeds per point; budget cap
    $200{,}000$ samples for sampling methods and $15$ seconds
    wall-clock for \toolname{}; budget-capped points are marked with
    hollow rings.  The dotted line is the Hoeffding envelope
    $\log(2/\delta) / (2\mutrue\varepsilon^2)$, the slope that plain
    MC must trace.  In the sweet-spot regime
    $\mutrue \in [5{\times}10^{-4}, 2{\times}10^{-3}]$,
    \code{dise\_betting} reaches half-width $0$ exactly (Y axis
    value~$\approx 10^3$, well below the MC envelope).}
  \label{fig:rare-event-scaling}
\end{figure}

\Cref{fig:rare-event-scaling} shows the result.  Three regimes are
visible.

\paragraph{Sweet spot, $\mutrue \in [5{\times}10^{-4},
  2{\times}10^{-3}]$: \toolname{} dominates.}
At $\code{rare\_width} = 20$ ($\mutrue \approx 0.002$),
\code{dise\_betting} certifies $\mutrue$ to half-width
\emph{exactly~$0$} in a median of $940$ samples across three seeds.
Plain Monte Carlo at the same target half-width reports
$\hat\mu = 0$ after $400$ samples---the rare event was never
observed, so the Wilson interval is tight around the wrong point
estimate.  Among the methods tested, \toolname{} is the only one
whose $\hat\mu$ tracks the true value $\mutrue = 0.002$ at small
sample budgets in this regime.

\paragraph{High-mass regime, $\mutrue \ge 10^{-2}$: plain MC
  dominates.}
Once $\mutrue$ is large enough that even a small sample sees several
hits, plain Monte Carlo's $\Theta(1/(\mutrue \varepsilon^2))$ scaling
beats \toolname{}'s per-iteration SMT cost.  At
$\code{rare\_width} = 1000$ ($\mutrue = 0.10$), plain Monte Carlo
reaches half-width $\varepsilon$ in $14{,}000$ samples while
\toolname{} hits the wall-clock cap.  The crossover sits at roughly
$\mutrue \approx 10^{-2}$ for this program structure.

\paragraph{Very rare regime, $\mutrue \le 10^{-4}$: both struggle.}
Plain Monte Carlo reports $\hat\mu = 0$---a tight but wrong interval.
\toolname{} hits the wall-clock cap because the $500$-sample
bootstrap is unlikely to observe the rare slice
($\mathbb{E}[\text{hits}] = 500 \cdot 10^{-4} = 0.05$); without an
observation in the rare branch, the SMT closure check on the dominant
\textsc{closed-f} leaf refuses to close.  A benchmark-aware
bootstrap policy would unblock this regime; we leave it to follow-up
work.

\subsection{Ablations}\label{sec:ablations}

\Cref{tab:ablations} reports the two ablations on
\code{assertion\_overflow\_mul\_w=8\_U(1,31)} (a non-LIA
multiplication benchmark, $\mutrue = 0.401$, budget $= 2000$).

\begin{table}[t]
  \centering
  \caption{Ablations on \code{assertion\_overflow\_mul\_w=8\_U(1,31)}
    at budget $= 2000$, three seeds, $\delta = 0.05$.  Median half-
    width with the median samples actually used.  Ablation~A: same
    bound on both sides, so the gap is entirely due to symbolic
    stratification.  Ablation~B: same adaptive-allocation strategy on
    both sides, so the gap is entirely due to symbolic guidance.}
  \label{tab:ablations}
  \small
  \begin{tabular}{lrrl}
    \toprule
    Method & Half-width & Samples & Note \\
    \midrule
    \multicolumn{4}{l}{\emph{Ablation~A: bound contribution vs.\ stratification contribution}} \\
    \code{plain\_mc\_betting}      & $0.0463$ & $2{,}000$ & PrPl-EB, no stratification \\
    \code{dise\_betting}           & $\mathbf{0.0307}$ & $\mathbf{220}$ & PrPl-EB, symbolic strata \\
    \midrule
    \multicolumn{4}{l}{\emph{Ablation~B: variance-driven vs.\ symbolic stratification}} \\
    \code{adaptive\_stratified}    & $0.1241$ & $2{,}000$ & Neyman + hash buckets \\
    \code{dise\_betting}           & $\mathbf{0.0307}$ & $\mathbf{220}$ & SMT path-condition strata \\
    \midrule
    \multicolumn{4}{l}{\emph{Reference: Wilson tightness on plain MC}} \\
    \code{plain\_mc}               & $0.0215$ & $2{,}000$ & Wilson is Bernoulli-tight \\
    \bottomrule
  \end{tabular}
\end{table}

Ablation~A shows that fixing the bound and varying the stratification
yields a $1.5\times$ tighter interval and---more strikingly---a
$9\times$ reduction in samples consumed.  Both methods carry the same
PrPl-EB bound at $\delta = 0.05$; the entire difference is
attributable to \toolname{}'s symbolic partition of the input domain
into three certifiable regions, two of which close via SMT after the
bootstrap.

Ablation~B isolates the value of \emph{symbolic} guidance over
\emph{variance-adaptive} sampling.  \code{adaptive\_stratified}
performs a two-pass Neyman allocation across $16$ random hash
strata---the strongest sampling-only stratifier we know.  Yet the
resulting interval is $4\times$ looser than \toolname{}'s; the
Bonferroni correction across $16$ buckets, each carrying its own
hit-rate uncertainty, washes out the variance reduction.
\toolname{}'s strata, in contrast, are SMT-certified
path-deterministic regions whose hit rate is known exactly after
closure.

The reference row shows that on the same benchmark, plain Wilson
\emph{at a fixed budget} gives a tighter interval than
\code{dise\_betting}---but consumes $9\times$ more samples.  The
right reading is that \toolname{} solves a different optimization:
it minimizes \emph{samples} for a target half-width, not half-width
for a fixed budget.

\paragraph{The bound ladder.}
\Cref{tab:bound-ladder} fixes the stratification (none) and varies
the per-sample bound construction on plain Monte Carlo.  Wilson is
tightest at finite $n$ for Bernoulli observations; Hoeffding is
distribution-free and loose; Maurer--Pontil empirical Bernstein is
variance-adaptive but loose at small $n$; WSR PrPl-EB pays an
anytime-validity premium that is recovered by the freedom to stop at
any sample count.

\begin{table}[t]
  \centering
  \caption{The bound ladder on \code{assertion\_overflow}: same
    estimator (plain MC), four different certified-interval
    constructions, budget $= 2000$, $\delta = 0.05$.}
  \label{tab:bound-ladder}
  \small
  \begin{tabular}{lll}
    \toprule
    Bound & Half-width & Validity regime \\
    \midrule
    Wilson (Bernoulli-tight) & $0.0215$ & fixed $n$ \\
    Hoeffding (textbook SMC) & $0.0304$ & fixed $n$, distribution-free \\
    Maurer--Pontil EB        & $0.0342$ & fixed $n$, variance-adaptive \\
    WSR PrPl-EB              & $0.0463$ & \textbf{anytime-valid}, variance-adaptive \\
    \bottomrule
  \end{tabular}
\end{table}

\subsection{Full Suite}\label{sec:suite}

\Cref{tab:suite} reports the full $12 \times 10$ grid at budget
$= 2000$.  Three patterns are immediate.

\paragraph{(i) Tautology benchmarks: plain MC wins, by design.}
\code{popcount\_w6}, \code{parity\_w6}, and \code{log2\_w6} have
$\mutrue = 1$ by construction; no conditional branch in the program
relates to the property; \toolname{} has nothing to refine on and
falls back to its open-leaf default.  Plain Monte Carlo certifies
$\hat\mu = 1$ in microseconds.  These benchmarks are kept in the
suite for completeness (they exercise the empty-path-clauses corner
of the algorithm) but should be read as \emph{outside \toolname{}'s
intended scope}.

\paragraph{(ii) Rare-event and assertion benchmarks: \toolname{}
matches or beats sampling-only.}
\code{coin\_machine} and \code{assertion\_overflow} are the
benchmarks where the program's branch structure aligns with the
input distribution.  Here \toolname{} certifies in $\le 1000$ samples
what plain MC needs $10^3$--$10^5$ samples to reach.

\paragraph{(iii) Non-LIA and BG-distributed benchmarks: \toolname{}
hits the wall-clock cap.}  On
\code{gcd\_steps\_le\_5\_BG}, \code{miller\_rabin\_w=2\_BG},
\code{collatz\_le\_30\_BG}, \code{modpow\_fits\_in\_4b}, and
\code{sieve\_primality}, the SMT cost of the closure check on
non-linear or BG-distributed path conditions saturates the $8$-second
cap before $\Wopen$ shrinks meaningfully.  The certified interval
collapses to $[0,1]$---sound but uninformative.  We discuss the
engineering paths to fix this in \Cref{sec:conclusion}.

\begin{table}[t]
  \centering
  \caption{Full suite, median half-width across three seeds at
    budget $= 2000$, $\delta = 0.05$.  Empirical coverage of the
    certified interval against the MC reference was $\ge 1 - \delta$
    on every \code{(benchmark, method)} cell except
    \code{quasi\_mc\_sobol} on the three BG benchmarks
    (\Cref{sec:discussion}).  Bold = best non-trivial half-width
    among the ten methods.}
  \label{tab:suite}
  \scriptsize
  \begin{tabular}{lrrrrr}
    \toprule
    Benchmark & \code{plain\_mc} & \code{pmc\_betting}
              & \code{adapt\_strat} & \code{dise\_wilson}
              & \code{dise\_betting} \\
    \midrule
    \code{coin\_machine\_U(1,9999)}      & $\mathbf{0.0041}$ & $0.0074$ & $0.0239$ & $0.0107$ & $0.0143$ \\
    \code{assertion\_overflow\_mul\_w=8} & $\mathbf{0.0215}$ & $0.0463$ & $0.1241$ & $0.0307$ & $0.0307$ \\
    \code{popcount\_w6}                  & $\mathbf{0.0005}$ & $0.0019$ & $0.0161$ & $0.500$  & $0.500$  \\
    \code{parity\_w6}                    & $\mathbf{0.0005}$ & $0.0019$ & $0.0161$ & $0.500$  & $0.500$  \\
    \code{log2\_w6}                      & $\mathbf{0.0005}$ & $0.0019$ & $0.0163$ & $0.500$  & $0.500$  \\
    \code{sparse\_trie\_depth\_le\_3}    & $\mathbf{0.0005}$ & $0.0019$ & $0.0163$ & $0.500$  & $0.500$  \\
    \code{integer\_sqrt\_correct}        & $\mathbf{0.0005}$ & $0.0019$ & $0.0163$ & $0.500$  & $0.500$  \\
    \code{sieve\_primality\_U(2,200)}    & $\mathbf{0.0185}$ & $0.0402$ & $0.1044$ & $0.500$  & $0.500$  \\
    \code{modpow\_fits\_in\_4b}          & $\mathbf{0.0218}$ & $0.0464$ & $0.1221$ & $0.500$  & $0.500$  \\
    \code{miller\_rabin\_w=2\_BG}        & $\mathbf{0.0209}$ & $0.0446$ & $0.0893$ & $0.500$  & $0.500$  \\
    \code{gcd\_steps\_le\_5\_BG}         & $\mathbf{0.0046}$ & $0.0078$ & $0.0251$ & $0.500$  & $0.500$  \\
    \code{collatz\_le\_30\_BG}           & $\mathbf{0.0112}$ & $0.0216$ & $0.0657$ & $0.500$  & $0.500$  \\
    \bottomrule
  \end{tabular}
\end{table}

The honest reading is that at fixed budget $= 2000$, plain Monte
Carlo wins on \emph{half-width} on most benchmarks.  This is the
\emph{wrong} question for \toolname{}, whose contract is
\emph{minimize samples for a target $\varepsilon$}, not minimize
$\varepsilon$ for a fixed sample count.
\Cref{fig:rare-event-scaling} reframes the question correctly and
exhibits the sweet-spot regime where \toolname{} delivers
$O(1)$-sample certification that plain MC cannot match at any budget.

\subsection{Discussion}\label{sec:discussion}

The empirical picture matches the theoretical analysis.  \toolname{}'s
structural variance reduction pays off when (a) the operational
distribution concentrates on path-deterministic regions of high
probability mass, (b) the bootstrap can observe at least one input
from each reachable region, and (c) the SMT cost of closure on the
program's path conditions is amortised over enough useful refinements.
The pedagogical \code{coin\_machine}, the assertion benchmark
\code{assertion\_overflow}, and the rare-slice parametric program of
\Cref{sec:rare-event-scaling} all satisfy (a)--(c).  On benchmarks
whose path conditions defeat Z3 within the per-cell budget---
multiplication, modular exponentiation, the GCD recurrence under a
BoundedGeometric distribution---\toolname{} cannot certify in
reasonable time and plain Monte Carlo is the right tool.

\paragraph{An honest negative: \code{quasi\_mc\_sobol} unsound on BG
benchmarks.}
On the three BoundedGeometric-distributed benchmarks
(\code{gcd\_steps\_le\_5\_BG}, \code{collatz\_le\_30\_BG},
\code{miller\_rabin\_w=2\_BG}), the \code{quasi\_mc\_sobol} baseline
reports $0\%$ empirical coverage.  The Wilson interval is constructed
assuming iid Bernoulli observations, but a scrambled Sobol sequence
is deterministic and structurally correlated.  On distributions whose
support is heavily concentrated on a small set of values---as a
BoundedGeometric with small $p$ is---the Sobol grid aligns with the
distribution in a way that bias-shifts the empirical mean outside the
Wilson interval.  This is a known property of quasi-MC paired with
naive intervals and is not a \toolname{} finding; it is included here
because the harness surfaces it honestly.

\paragraph{Threats to validity.}
The benchmark suite contains $12$ programs and the rare-event
parametric experiment varies one knob; broader empirical coverage
would strengthen the claim.  All measurements are on a single
hardware configuration with a single SMT backend (Z3); SMT cost in
particular is implementation-dependent.  Three seeds give a sense of
variance but do not yield narrow confidence intervals on the medians
reported.  We have not compared head-to-head against PSI~\cite{psi}
or Dice~\cite{holtzen2020dice}; both require porting our benchmarks
to dedicated input languages and are deferred to follow-up work.

\section{Conclusion}\label{sec:conclusion}

We have presented \algname{}, an algorithm for distribution-aware
reliability estimation that maintains a frontier of path-condition
regions and at each step chooses between sampling and SMT-driven
refinement under a unified gain-per-cost criterion.  The certified
half-width decomposes into a statistical component $\epsstat$, driven
down by sampling, and a structural component $\Wopen$, driven down by
symbolic refinement.  This decomposition---Theorem~\ref{thm:hw-decomp}---is
simultaneously the algorithm's correctness argument and its scheduler.
Two implementation choices distinguish \toolname{} from prior
work: the per-leaf statistical bound instantiates the closed-form
predictable-plug-in empirical-Bernstein confidence sequence of Waudby-Smith
and Ramdas, and the closure rule is single-branch (SMT-only), yielding
a one-lemma soundness contract and eliminating a class of silent bias
that affects prior prototypes.

\paragraph{Limitations.}
We identify five.

\begin{enumerate}
\item \emph{Path-determinism versus property-constancy.}
  \Cref{sec:soundness-caveat} reports a measured soundness regression
  on a non-LIA benchmark where the property is evaluated outside the
  concolic tracer.  The fix is well-understood but not yet
  implemented.

\item \emph{Distribution class.}  The (D1) restriction excludes
  Bayes-net-structured and general discrete distributions.  Lifting it
  requires ancestor-conditioned constrained samplers and a mass
  estimator that handles coordinate dependence.

\item \emph{Floating-point arithmetic and symbolic memory.}
  Both are out of scope for the current prototype.

\item \emph{Refinement-clause selection.}  The current heuristic
  picks from clauses observed in samples; it does not synthesise new
  predicates.  Craig interpolation is the natural principled
  alternative.

\item \emph{Statistical efficiency under partition fragmentation.}
  When the frontier fragments into many high-mass open leaves, the
  Bonferroni correction inflates $\epsstat$ in proportion to the leaf
  count, and the SMT cost of refinement is not recovered by variance
  reduction.
\end{enumerate}

\paragraph{Future work.}
Two extensions deserve immediate attention.  First, the closure rule's
implicit path-implies-property assumption can be discharged by routing
$\prop$ through the concolic tracer---straightforward, requiring only
careful handling of short-circuit evaluation.  Second, replacing the
candidate-clause scorer with Craig
interpolation~\cite{mcmillan2003interpolation,henzinger2004abstraction}
would let the algorithm synthesise predicates the program never
explicitly branches on, broadening the class of programs for which
$\Wopen$ can be driven to zero.  Beyond these, we are interested in
substituting a certified approximate model
counter~\cite{kiesl2024certified,yang2023approxmc6} for the
importance-sampling estimator on non-axis-aligned regions---yielding
machine-checkable certificates end-to-end---and in extending the
distribution class to handle (D2) Bayes-net-structured inputs.

\paragraph{Reproducibility.}
\toolname{}, including all benchmarks, scripts, and the Docker image,
is anonymised for review at \code{https://github.com/<anon>/dise}.
All tables and figures in this paper are reproduced end-to-end by
\code{scripts/reproduce.sh}.


% ---------------------------------------------------------------------------
% Bibliography
% ---------------------------------------------------------------------------
\bibliographystyle{splncs04}
\bibliography{references}

\end{document}
